\section{Party restrictions after changing owner labels}

Let $P$ and $Q$ be sets of owner labels. A finite composition $W$ contains
local operations, exchanges, and operations beside untouched registers.
Write $\operatorname{Part}(W)$ for its participating parties: these include the owners of
all occurring registers and the assigned party of every local operation,
including a scalar operation on an empty register list. For $f:P\to Q$,
write $W^f$ for the composition obtained by changing each owner $p$ to
$f(p)$ while preserving the register spaces and chronological order.

\begin{definition}[Selection of one party]
    \label{def:peps_single_party_selector}
    \lean{TNLean.PEPS.PairEffect.partySelector}
    \leanok
    For $q\in Q$, define $\chi_q:Q\to\{\mathrm{false},\mathrm{true}\}$
    by $\chi_q(x)=\mathrm{true}$ exactly when $x=q$.
\end{definition}

\begin{theorem}[Participating parties after changing owner labels]
    \label{thm:peps_source_free_owner_parties}
    \lean{TNLean.PEPS.PairEffect.Word.parties_mapOwner_of_sources_nil}
    \leanok
    \uses{def:peps_owner_composition}
    If $W$ contains no pair-source preparation, then
    $\operatorname{Part}(W^f)=f(\operatorname{Part}(W))$. No injectivity assumption on $f$ is required.
\end{theorem}

\begin{proof}\leanok
    \uses{def:peps_owner_composition}
    Induct on the composition. The owners of a relabelled register list
    are the image of its original owners. Images commute with unions
    for composition and with adjoining a named party for local maps,
    exchanges, and untouched registers. The pair-source case is excluded
    by hypothesis. This also retains named scalar local operations.
\end{proof}

\begin{theorem}[Restriction to an owner outside the relabelling range]
    \label{thm:peps_owner_disjoint_restriction}
    \lean{TNLean.PEPS.PairEffect.Word.restrict_mapOwner_eval_eq_id_of_disjoint}
    \leanok
    \uses{def:peps_single_party_selector,def:peps_owner_composition}
    Suppose that $W$ contains no pair-source preparation and that
    $f(p)\ne q$ for every $p\in P$. Let $V$ be the restriction of $W^f$
    to the owner $q$. Then $\operatorname{Part}(V)=\varnothing$, and its operator is
    $\operatorname{id}_{\C}$ under the canonical identification of the
    empty input and output register memories with $\C$.
\end{theorem}

\begin{proof}\leanok
    \uses{thm:peps_source_free_owner_parties,thm:peps_party_restriction,
        def:peps_single_party_selector,thm:peps_party_factorization}
    The participating parties are
    $\operatorname{Part}(V)=\{x\in f(\operatorname{Part}(W)):x=q\}=\varnothing$.
    A composition with no participating party has empty input and output
    register lists and acts as their identity. In particular, it cannot
    contain a scalar local operation with a named party.
\end{proof}

% Source: 04-compression.tex, lines 233--267 and 351--381.
These auxiliary identities concern the separated local operations in
\cite[proof of Theorem~5.2]{OpenAI2026PolynomialPEPS}.
For the all-party private preparation, every original party has its own
owner in $P\sqcup\{\bot\}$. Its restriction to $\bot$ is therefore the
scalar identity. Adjoining fixed source registers requires, in addition,
that none of those registers has exterior owner $\bot$. The subsequent
residual composition can contribute a separate scalar from operations on
empty memories; that scalar is not determined by the preparation identity.
